% GENERATED FILE. Edit canonical lessons, then run npm run generate:books.
% canonical-source-hash: fnv1a64:401dfc290fb1d9c4
% canonical-lessons: LA-C357-arbitror, LA-C357-existimo, LA-C357-censeo, LA-C357-opinor, LA-C357-opinio

\chapter{What I Think and Judge}
\label{ch:la-what-i-think-and-judge}
\hlchaptermodality{357}

\begin{chapteropening}
\textbf{By the end of this chapter:} \emph{I can say what I judge, reckon, propose and suppose, and talk about an opinion.}

Complete the last lesson of chapter 357: I can say what I judge, reckon, propose and suppose, and talk about an opinion.
\end{chapteropening}

\section[arbitror, arbitrārī]{arbitror, arbitrārī — to be of the opinion, to judge}

\label{lesson:LA-C357-arbitror}



\begin{quote}
Before the new one: say the Latin for a ditch, then the Latin for a rampart.
\end{quote}

\subsection*{You'll want to know: arbitror, arbitrārī}
\textbf{arbitror, arbitrārī} — \textquotedblleft{}to be of the opinion, to judge\textquotedblright{}.

\begin{quote}
Arbitror eum vērum dīcere.
\end{quote}

\begin{quote}
\emph{I believe he is telling the truth.}
\end{quote}

Deponent: \emph{arbitror, arbitrārī, arbitrātus sum}, passive in form and active in meaning. Like \emph{putō} it takes the accusative and infinitive. English \emph{arbitrate} comes from it.

\begin{cousinweb}
\emph{arbitror} is built on \emph{arbiter} \textquotedblleft{}an eyewitness, an umpire\textquotedblright{} (Lewis \& Short; EtymDB): to judge a thing as one who saw it. English \emph{arbiter}, \emph{arbitrate} and \emph{arbitrary} come from the same family.
\end{cousinweb}

\begin{grammarlens}[title={What you've built}]
One more word for telling, explaining and describing at B1.
\end{grammarlens}

\subsection*{Your turn}
\begin{itemize}
\raggedright
  \item \emph{Say it:} \emph{arbitror, arbitrārī}
  \item \emph{Say it:} \emph{arbitror, arbitrārī}, once more
  \item \emph{Say it:} a sentence of your own with \emph{arbitror, arbitrārī}
  \item \emph{Recall:} say the Latin for a ditch, then the Latin for a rampart, then say \emph{arbitror, arbitrārī} again
\end{itemize}

\begin{tcolorbox}[breakable,colback=teal!4,colframe=teal!35!black,title={Before you move on}]
What does arbitror, arbitrārī mean? (\textquotedblleft{}To be of the opinion, to judge\textquotedblright{}.) Say it once more.
\end{tcolorbox}
\section[exīstimō, exīstimāre]{exīstimō, exīstimāre — to reckon, to consider; to value}

\label{lesson:LA-C357-existimo}



\begin{quote}
Before the new one: say the Latin for a rampart, then the Latin for to be of the opinion.
\end{quote}

\subsection*{You'll want to know: exīstimō, exīstimāre}
\textbf{exīstimō, exīstimāre} — \textquotedblleft{}to reckon, to consider; to value\textquotedblright{}.

\begin{quote}
Quid dē hōc librō exīstimās?
\end{quote}

\begin{quote}
\emph{What do you think of this book?}
\end{quote}

\emph{exīstimō, exīstimāre, exīstimāvī, exīstimātum}. With the accusative and infinitive: \emph{exīstimō tē errāre}, I reckon you are wrong.

\begin{cousinweb}
\emph{exīstimō} is built from \emph{ex} \textquotedblleft{}out\textquotedblright{} and \emph{aestimō} \textquotedblleft{}to value, to put a price on\textquotedblright{} (Lewis \& Short): to weigh a thing up and give the result. English \emph{estimate} and \emph{esteem} come from \emph{aestimō}.
\end{cousinweb}

\begin{grammarlens}[title={What you've built}]
One more word for telling, explaining and describing at B1.
\end{grammarlens}

\subsection*{Your turn}
\begin{itemize}
\raggedright
  \item \emph{Say it:} \emph{exīstimō, exīstimāre}
  \item \emph{Say it:} \emph{exīstimō, exīstimāre}, once more
  \item \emph{Say it:} a sentence of your own with \emph{exīstimō, exīstimāre}
  \item \emph{Recall:} say the Latin for a rampart, then the Latin for to be of the opinion, then say \emph{exīstimō, exīstimāre} again
\end{itemize}

\begin{tcolorbox}[breakable,colback=teal!4,colframe=teal!35!black,title={Before you move on}]
What does exīstimō, exīstimāre mean? (\textquotedblleft{}To reckon, to consider; to value\textquotedblright{}.) Say it once more.
\end{tcolorbox}
\section[cēnseō, cēnsēre]{cēnseō, cēnsēre — to give an opinion, to propose; to assess}

\label{lesson:LA-C357-censeo}



\begin{quote}
Before the new one: say the Latin for to be of the opinion, then the Latin for to reckon.
\end{quote}

\subsection*{You'll want to know: cēnseō, cēnsēre}
\textbf{cēnseō, cēnsēre} — \textquotedblleft{}to give an opinion, to propose; to assess\textquotedblright{}.

\begin{quote}
Ego cēnseō pācem meliōrem esse quam bellum.
\end{quote}

\begin{quote}
\emph{In my opinion peace is better than war.}
\end{quote}

\emph{cēnseō, cēnsēre, cēnsuī, cēnsum}. The verb of a senator stating his view in the senate, and of the officer who counted and valued the citizens. English \emph{census} comes from it.

\begin{cousinweb}
\emph{cēnseō} is inherited from Proto-Italic \emph{*kensēō} \textquotedblleft{}to assess, to judge\textquotedblright{} (EtymDB), and has a Sanskrit cousin meaning \textquotedblleft{}to declare\textquotedblright{}. English \emph{census} comes from it.
\end{cousinweb}

\begin{grammarlens}[title={What you've built}]
One more word for telling, explaining and describing at B1.
\end{grammarlens}

\subsection*{Your turn}
\begin{itemize}
\raggedright
  \item \emph{Say it:} \emph{cēnseō, cēnsēre}
  \item \emph{Say it:} \emph{cēnseō, cēnsēre}, once more
  \item \emph{Say it:} a sentence of your own with \emph{cēnseō, cēnsēre}
  \item \emph{Recall:} say the Latin for to be of the opinion, then the Latin for to reckon, then say \emph{cēnseō, cēnsēre} again
\end{itemize}

\begin{tcolorbox}[breakable,colback=teal!4,colframe=teal!35!black,title={Before you move on}]
What does cēnseō, cēnsēre mean? (\textquotedblleft{}To give an opinion, to propose; to assess\textquotedblright{}.) Say it once more.
\end{tcolorbox}
\section[opīnor, opīnārī]{opīnor, opīnārī — to suppose, to imagine}

\label{lesson:LA-C357-opinor}



\begin{quote}
Before the new one: say the Latin for to reckon, then the Latin for to give an opinion.
\end{quote}

\subsection*{You'll want to know: opīnor, opīnārī}
\textbf{opīnor, opīnārī} — \textquotedblleft{}to suppose, to imagine\textquotedblright{}.

\begin{quote}
Ille, opīnor, nōn veniet.
\end{quote}

\begin{quote}
\emph{He will not come, I suppose.}
\end{quote}

Deponent: \emph{opīnor, opīnārī, opīnātus sum}. Dropped into a sentence, \emph{opīnor} softens it, like English \emph{I suppose}. English \emph{opine} comes from it.

\begin{cousinweb}
\emph{opīnor} comes from an old adjective \emph{opīnus} \textquotedblleft{}expecting\textquotedblright{} (EtymDB), which survives in \emph{necopīnus} \textquotedblleft{}unexpected\textquotedblright{} (Lewis \& Short). English \emph{opine} and \emph{opinion} come from it.
\end{cousinweb}

\begin{grammarlens}[title={What you've built}]
One more word for telling, explaining and describing at B1.
\end{grammarlens}

\subsection*{Your turn}
\begin{itemize}
\raggedright
  \item \emph{Say it:} \emph{opīnor, opīnārī}
  \item \emph{Say it:} \emph{opīnor, opīnārī}, once more
  \item \emph{Say it:} a sentence of your own with \emph{opīnor, opīnārī}
  \item \emph{Recall:} say the Latin for to reckon, then the Latin for to give an opinion, then say \emph{opīnor, opīnārī} again
\end{itemize}

\begin{tcolorbox}[breakable,colback=teal!4,colframe=teal!35!black,title={Before you move on}]
What does opīnor, opīnārī mean? (\textquotedblleft{}To suppose, to imagine\textquotedblright{}.) Say it once more.
\end{tcolorbox}
\section[opīniō, opīniōnis]{opīniō, opīniōnis — an opinion, a belief; a reputation}

\label{lesson:LA-C357-opinio}



\begin{quote}
Before the new one: say the Latin for to give an opinion, then the Latin for to suppose.
\end{quote}

\subsection*{You'll want to know: opīniō, opīniōnis}
\textbf{opīniō, opīniōnis} — \textquotedblleft{}an opinion, a belief; a reputation\textquotedblright{}.

\begin{quote}
Haec est opīniō mea, nōn tua.
\end{quote}

\begin{quote}
\emph{This is my opinion, not yours.}
\end{quote}

\emph{opīniō, opīniōnis} f., third declension. It is also what people think of you: \emph{bona opīniō}, a good name. English \emph{opinion}.

\begin{cousinweb}
\emph{opīniō} is the noun built on \emph{opīnor} \textquotedblleft{}to suppose\textquotedblright{} (EtymDB; Lewis \& Short), and English \emph{opinion} comes from it through French.
\end{cousinweb}

\begin{grammarlens}[title={What you've built}]
One more word for telling, explaining and describing at B1.
\end{grammarlens}

\subsection*{Your turn}
\begin{itemize}
\raggedright
  \item \emph{Say it:} \emph{opīniō, opīniōnis}
  \item \emph{Say it:} \emph{opīniō, opīniōnis}, once more
  \item \emph{Say it:} a sentence of your own with \emph{opīniō, opīniōnis}
  \item \emph{Recall:} say the Latin for to give an opinion, then the Latin for to suppose, then say \emph{opīniō, opīniōnis} again
\end{itemize}

\begin{tcolorbox}[breakable,colback=teal!4,colframe=teal!35!black,title={Before you move on}]
What does opīniō, opīniōnis mean? (\textquotedblleft{}An opinion, a belief; a reputation\textquotedblright{}.) Say it once more.
\end{tcolorbox}
